\section{The Pair Mutual-Correlated-Agreement Theorem}

\subsection{Bad challenges}

Let \(p=(p_0,p_1)\in(\F^2)^I\), and write
\[
  y_\gamma=p_0+\gamma p_1.
\]
For \(S\subseteq I\), say that \(S\) is \(\delta\)-large when
\[
  \frac{|S|}{n}\ge1-\delta.
\]
Define \(\Bad_{\MCA}(\mathcal C,\delta,p)\) to be the set of
\(\gamma\in\F\) for which there exists a \(\delta\)-large set \(S\) such that:
\begin{enumerate}[label=(\roman*)]
  \item \(y_\gamma\) agrees on \(S\) with a codeword of \(\mathcal C\);
  \item the pair \(p\) does not agree on the same \(S\) with a codeword of
    \(\Cpair\).
\end{enumerate}
The MCA error is
\[
  \varepsilon_{\MCA}(\mathcal C,\delta)
  =
  \max_{p\in(\F^2)^I}
  \frac{|\Bad_{\MCA}(\mathcal C,\delta,p)|}{|\F|}.
\]

\begin{theorem}[Uniform pair-MCA certificate bound]
\label{thm:pair-mca}
Let \(\mathcal C=\RS[\F,\mathcal L,k]\) have injectively indexed
evaluation domain of length \(n\), where \(1\le k\le n\).  For every
\(\delta\in\mathbb R\) and every \(p\in(\F^2)^n\),
\[
  |\Bad_{\MCA}(\mathcal C,\delta,p)|
  \le \binom n{k+1}.
\]
Consequently,
\[
  \varepsilon_{\MCA}(\mathcal C,\delta)
  \le \frac{\binom n{k+1}}{|\F|}.
\]
Smoothness is not required for this bound.
\end{theorem}

\begin{proof}
Fix a bad challenge \(\gamma\), and choose a witnessing set \(S\).
Necessarily \(|S|>k\).  Indeed, when \(|S|\le k\), interpolate each component
of \(p\) on \(S\) by a polynomial of degree \(<k\); the resulting pair
codeword agrees with \(p\) on \(S\), contradicting condition~(ii).

Choose a \(k\)-subset \(T\subset S\).  Let \(a_T,b_T\) be the unique
degree-\(<k\) polynomials interpolating \(p_0,p_1\), respectively, on \(T\).
Condition~(ii) says that the pair interpolant \((a_T,b_T)\) fails somewhere
on \(S\). Choose such an \(i\in S\setminus T\), and set
\(U=T\cup\{i\}\). Condition~(i) remains true after restriction from \(S\) to
\(U\). The pair still cannot agree with \(\Cpair\) on \(U\): any purported
pair codeword would agree with \((a_T,b_T)\) on \(T\), hence equal it by
\(k\)-point determination, contradicting the mismatch at \(i\).
Thus \(U\) is a \((k+1)\)-coordinate obstruction for \(\gamma\).

A fixed \(U\) obstructs at most one challenge. Suppose distinct
\(\gamma,\eta\) were both bad on \(U\), with scalar codewords
\(c_\gamma,c_\eta\) agreeing with
\(p_0+\gamma p_1,p_0+\eta p_1\), respectively. By linearity,
\[
  d=\frac{c_\gamma-c_\eta}{\gamma-\eta},
  \qquad
  a=c_\gamma-\gamma d
\]
are codewords. On \(U\), they agree with \(p_1\) and \(p_0\), respectively,
so \((a,d)\) is a pair codeword agreeing with \(p\) on \(U\), a
contradiction.

Make the finite choices deterministic. Mapping each bad challenge to its
unordered obstruction \(U\) is injective, and there are
\(\binom n{k+1}\) such subsets.
Dividing by \(|\F|\) proves the error bound.  Notice that the large-set
inequality is no longer used after the witness is extracted; this is why the
bound is uniform in \(\delta\).
\end{proof}

\begin{corollary}[Explicit smooth family at \(2^{-128}\)]
\label{cor:explicit-family-mca}
For every member \(\mathcal C_u\) of the family in
\cref{sec:explicit-family}, every \(\delta\in[0,1]\), and every pair word
\(p\),
\[
  \frac{|\Bad_{\MCA}(\mathcal C_u,\delta,p)|}{|\F_u|}
  \le 2^{-128}.
\]
Thus, for this family and this same-set pair-MCA definition, every normalized
radius is safe and the largest admissible radius is \(1\).
\end{corollary}

\begin{proof}
\Cref{sec:field-inequality} proves
\[
  2^{128}\binom{n_u}{k_u+1}\le q_u.
\]
Apply \cref{thm:pair-mca}.
\end{proof}

\subsection{An exact small-radius boundary instance}
\label{sec:mca-boundary-field}

The uniform certificate bound is sufficient but need not locate the first
unsafe radius.  Write
\[
  B_{\mathcal C}(r)
  =
  \max_p|\Bad_{\MCA}(\mathcal C,r/n,p)|
\]
for the exact worst bad-challenge count at count radius \(r\).  Consider
\[
  q=2^{128}+192513
   =340282366920938463463374607431768403969.
\]
The accompanying Lean development certifies that \(q\) is prime and
\(q\equiv1\pmod {4096}\).  Hence \(\mathbb F_q^\times\) contains an
order-\(4096\) subgroup, which we use as the evaluation domain for the smooth
RS code with \(n=4096\) and \(k=1024\).

\begin{theorem}[Exact 128-bit small-radius transition]
\label{thm:mca-boundary-threshold}
For this code,
\[
  B_{\mathcal C}(0)=1,
  \qquad
  B_{\mathcal C}(1)\ge2.
\]
Consequently its greatest safe count radius at 128 bits is \(0\), and its
first unsafe count radius is \(1\).  For the relative-radius definition over
all \(\delta\in[0,1]\), the safe set is
\[
  \left[0,\frac1{4096}\right),
\]
so the real-radius supremum is \(1/4096\) and is not attained.
\end{theorem}

\begin{proof}
At count radius zero, every large agreement set is the whole coordinate set.
A fixed set supports at most one bad challenge, while a unit-spike affine
line realizes one; hence \(B_{\mathcal C}(0)=1\).  At count radius one, take
two coordinates outside a \(k\)-point interpolation set.  The affine line
from the unit spike at the first coordinate to the unit spike at the second
has distinct bad challenges \(0\) and \(1\), proving
\(B_{\mathcal C}(1)\ge2\).

The certified field lies in the exact budget window
\[
  2^{128}\le q<2^{128}\cdot2.
\]
Thus the integer criterion \(2^{128}B_{\mathcal C}(r)\le q\) holds at
\(r=0\) and fails at \(r=1\).  Monotonicity of the bad-count spectrum gives
the two grid endpoints.

Finally, if \(0\le\delta<1/n\), the inequality
\(|S|/n\ge1-\delta>1-1/n\) forces the integer \(|S|\) to equal \(n\).
The bad sets are therefore exactly those at count radius zero throughout
that half-open interval.  At \(\delta=1/n\), sets missing one coordinate
become admissible, and the two-spike line already violates the budget.
Monotonicity preserves failure for larger radii, proving the real-radius
claim.
\end{proof}

\subsection{Distance interpretation}

Let
\[
  D_2(p)=\Delta(p,\Cpair),
  \qquad
  D_1(y)=\Delta(y,\mathcal C).
\]
For a challenge outside the MCA bad set, scalar closeness on a large agreement
set forces pair closeness on that same set.  In distance language, the desired
contrapositive has the form
\[
  D_2(p)\ge\delta
  \Longrightarrow
  D_1(y_\gamma)\ge\delta
\]
for every MCA-good \(\gamma\), with the fraction of excluded challenges bounded
by \(\varepsilon_{\MCA}(\mathcal C,\delta)\).

This is a challenge-indexed statement.  Replacing the sampled challenge by one
fixed field element would require a separate theorem that the selected element
is good for every pair, which does not follow from a fractional MCA bound.

\subsection{Endpoint interpretation}

At \(\deltacap=1-k/n\), interpolation on any \(k\) coordinates shows that
\[
  D_1(y)\le\deltacap
  \quad\text{and}\quad
  D_2(p)\le\deltacap
\]
for every scalar word \(y\) and pair word \(p\).  Moreover, on every set
\(S\) of size at most \(k\), each component of \(p\) extends to a
degree-\(<k\) polynomial.  Such an \(S\) can therefore never witness an MCA
failure.  Every bad endpoint challenge must instead have a witness with
\(|S|\ge k+1\).

Consequently, the endpoint distance implication above is an equality
non-collapse statement: when \(D_2(p)=\deltacap\), a good challenge must not
lower the projected distance strictly below \(\deltacap\).
\Cref{thm:pair-mca} supplies the previously missing quantitative bound on the
excluded challenges.

\subsection{Commit-then-challenge order}

The mathematical experiment fixes \(p\) before sampling
\(\gamma\leftarrow\F\).  This ordering prevents the pair from being selected as
a function of the challenge.  A full protocol implementation must bind \(p\)
before deriving \(\gamma\), but the coding theorem needs only the resulting
quantification and probability space; it does not depend on a complete
blockchain or commitment protocol.

\subsection{Scope}

The theorem is universal over pair words and injective RS evaluation domains,
but its \(2^{-128}\) corollary uses the explicit exponential-size field.
For a different supplied field, the same certificate proof gives the
sufficient comparison
\[
  \binom n{k+1}\le 2^{-128}|\F|.
\]
This improves the ordered-certificate requirement by the exact factor
\(k+1\), since
\(\binom nk(n-k)=(k+1)\binom n{k+1}\).
If that inequality fails, this proof does not determine the optimal radius or
show that the MCA target is impossible; it only ceases to certify the requested
error.

For an arbitrary supplied finite code, radius \(r/n\) meets the
\(2^{-128}\) target exactly when
\[
  2^{128}B_{\mathcal C}(r)\le|\F|.
\]
This is also the exact formal equivalence: the Lean theorem identifies the
real-valued operational error predicate at normalized radius \(r/n\) with the
displayed integer inequality, so no rounding relaxation is hidden here.
Thus a uniform upper-bound certificate at \(r\), together with one concrete
affine line exceeding the budget at \(r+1\), pins both the largest safe and
first unsafe Hamming-grid radii. The remaining middle regime is the
code-specific problem of proving those matching bounds; smoothness by itself
does not determine this incidence spectrum.

The result is for the two-coordinate, same-set MCA experiment defined above.
Broader notions or higher-dimensional challenge families require
separate statements.

As a separate concrete/application result, the repository also checks the
canonical runtime field \(q=12289\) with \(n=4096\) and \(k=1024\).  Since
\(q<2^{128}\) while its radius-zero unit-spike line has one bad challenge,
that concrete code has no safe count radius at 128-bit MCA security.  This
runtime-facing fact is not used by the standalone large-field or boundary-field
theorems above.
